% Transcription — fonds Alexandre Grothendieck, shelfmark 146, batch 3 (pages 41-60).
% Established from the facsimile archives/batches/146/batch-03.pdf (a mirror of
% https://grothendieck.umontpellier.fr/146.pdf), per the skill
% transcribe-grothendieck.
%
% Pass: Opus 5.5 (claude-opus-5-5), 2026-10-02 — first pass, unchecked against the pages by a human.
%
% Fourteen of the twenty pages are transcribed. Skipped: 56 (blank), 58 and
% 60 (USTL thesis-defence notices of 1979, typed, nothing in his hand).
% Page 41 is a very faint ghost of page 42's text, line for line; it gets
% one \note and is not transcribed twice. Page 55 carries only a pencil
% heading.
%
% Pages 42-54 — a fair-ish draft, his pages 1-7: a regular 2n-gon in a real
%   plane V, its dihedral group D = Aut(V, J~), orientations ω, the central
%   symmetry a_0, the reflections r_j; then the complexification V_C,
%   Σ = P(V_C), the two poles P_ω, the isotropic splitting V_C = ⊕ V_C^ω,
%   the bundle O_Σ(1) and its C*-torsor PΣ ≃ V*, the involution σ, and the
%   open sets Σ* = Σ ∖ {R_i}, PΣ* ≃ V** whose fundamental groupoids he wants
%   by generators and relations, with the action of D × μ_m × {1, τ}. The
%   NB notes written slantwise in the margins (pp. 44, 45, 50, 51) are
%   largely illegible.
% Pages 55-59 — headed "0T05 et S0T0,5": sketches of spheres with marked
%   circles and points, counts of "possibilités pour K" (60, 12, 20, 30),
%   types I-IV, D_0, D_1, D_2, D_01, D_02 and the icosahedron.
%
% Equation numbers are his, and he repeats (18), (19), (21).
%
\documentclass[11pt,a4paper]{article}
\input{../preamble/grothendieck}

\folder{146}
\batch{3}
\pages{41}{60}
\foldertitle{Groupoïdes de Teichmüller et les S0 Tg,!ν, S0 TG : notes manuscrites (s.d.).}
\dating{[à partir de 1978]}
\watermark{Édition de démonstration}

\begin{document}

\page{41}
\note{la feuille ne porte, très pâle, que le texte de la p.~42,
ligne pour ligne et note marginale comprise (« Soient $n$ un entier
$\geq 3$, et considérons un polygone régulier… », jusqu'à
$z \mapsto \zeta z$) ; il est transcrit une seule fois, à la page
suivante.}

\page{42}
\marginal{Il suffirait \ill{} $n \geq 2$, mais le cas $n = 2$ est un
\uncertain{peu} très dégénéré \ill{} \uncertain{pour certaines} \ill{} ;
\ill{} manque d'ailleurs \ill{} distinct\ill{}, \ill{} c'est le cas du
\ill{} \uncertain{un trou} !}
Soient $n$ un entier $\geq 3$, et considérons un polygone régulier à
\add{$m =$} $2n$ côtés, dans un espace vectoriel \add{réel} $V$ de
dimension 2. L'ens.\ $\tilde{J}$ des sommets du polygone forme
\add{l'ens.\ des sommets d'}un polygone combinatoire (on dira aussi qu'il
est muni d'une structure de polygone combinatoire, exprimable de diverses
façons équivalentes bien connues), dont la donnée équivaut, à isom.\ can.\
près, à celle du plan vectoriel réel $V$ et du polygone régulier dans $V$.
La donnée de l'un ou de l'autre équivaut en effet à celle d'un torseur
sous le groupe diédral
\[
  (1)\qquad D_m = \mu_m \cdot \mathbb{Z}/2\mathbb{Z}
\]
(produit semi-direct, $\mathbb{Z}/2\mathbb{Z}$ opérant sur $\mu_m$ via
$\zeta \mapsto \zeta^{-1}$) qui est le groupe d'automorphismes de la
structure standard
\[
  (2)\qquad V_0 = \mathbb{C}, \qquad \tilde{J}_0 = \mu_m
  \overset{\text{déf}}{=} {}_m\mathbb{C}^{*} \subset V_0 ,
\]
$\mu_m$ opérant sur $(V_0, \tilde{J}_0)$ par homothéties
\[
  z \longmapsto \zeta z \qquad (\zeta \in \mu_m)
\]
et $\mathbb{Z}/2\mathbb{Z}$ par $z \longmapsto \bar{z}$ conjugaison complexe.

\page{43}
On posera
\[
  (3)\qquad D = D_{\tilde{J}} = \operatorname{Aut}(V, \tilde{J})
  \qquad (\tilde{J} \subset V)
\]
\[
  (4)\qquad \underline{\omega} = \underline{\omega}_{\tilde{J}}
  = \Omega(\tilde{J}) \simeq \Omega(V)
\]
\marginal{ens.\ des orientations du poly.\ comb.\ $\tilde{J}$, ou de
l'espace vectoriel réel $V$}
On a dans $D$ le sous-groupe $D^{+}$ des rotations, qui est le noyau de
l'hom.\ naturel \add{(transport de str.)}
$D \to \mathbb{Z}/2\mathbb{Z} \simeq \operatorname{Aut}(\underline{\omega})$,
d'où suite exacte
\[
  (5)\qquad 1 \longrightarrow D^{+} \longrightarrow D \longrightarrow
  \mathbb{Z}/2\mathbb{Z} \longrightarrow 1 ,
\]
$\mathbb{Z}/2\mathbb{Z}$ opérant sur $D^{+}$ par $\zeta \mapsto \zeta^{-1}$.
\struck{On a un plongement \ill{}}
\[
  (6)\qquad \underline{\omega} \subset D^{+}
\]
et $\underline{\omega}$ est formé de deux rotations inverses l'une de
l'autre, d'angle $\frac{2\pi}{m}$ \struck{pour l'} \add{dans le} sens de
\struck{l'orientation} \add{l'une et l'autre} \uncertain{orientation} de
$V$. Ce sont deux générateurs de $D^{+}$, inverses \struck{l'un de l'autre.
On a si} l'un de l'autre, $\omega' = \omega^{-1}$, iso correspondant
(\uncertain{le choix de} $\omega$ \ill{} $\mathbb{Z}/m\mathbb{Z}
\xrightarrow{\sim} D^{+}$). \struck{Les points de} Les sommets du polygone
seront désignés individuellement par $\tilde{R}_i$ ($i \in \tilde{J}$).
On n'a pas utilisé \struck{\ill{}} le fait que $m$ est pair, ce qui
équivaut au fait que l'ens.\ des $\tilde{R}_i$ dans $V$ est stable par la
symétrie $x \mapsto -x$,

\page{44}
\note{p.~2 de l'auteur.}
\struck{donc} \add{on a encore} \uncertain{défini} par un élément \struck{privilégié
d'ordre} \add{d'ordre 2}
\[
  (7)\qquad a_0 \in D^{+}, \qquad a_0 : x \longmapsto -x ,
\]
\struck{qui est} en fait le seul élément d'ordre 2 de $D^{+}$. Si à tout
$j \in \tilde{J}$ on associe l'unique \struck{réflexion} élément
$r_j \in D$ (\uncertain{nécessairement une} réflexion) \struck{donc} qui soit
$\neq 1$ et telle que
\[
  (8)\qquad r_j(\tilde{R}_j) = \tilde{R}_j ,
\]
on \uncertain{trouve}
\[
  (9)\qquad r_j = r_{a_0(j)} ,
\]
\struck{et} donc il y a \uncertain{précisément} $n = m/2$ réflexions $r_j$
distinctes, qui correspondent à la moitié \add{\uncertain{exactement}} des
éléments de la classe $D^{-}$ de $D$ modulo $D^{+}$. Les autres réflexions
sont \uncertain{celles} associées aux \emph{côtés} des polygones envisagés.

Il y a sur $V$ une unique forme quadratique invariante par $D$ (\ill{}
seulement par $D^{+}$), \uncertain{prise} égale à 1 sur les $\tilde{R}_j$.
On considère $V$ comme espace vect.\ euclidien, grâce à cette forme.

\marginal{NB Le groupe \struck{agissant} $D$ est \uncertain{réduit} à
$\{1, a_0\}$ \ill{} \ill{} cas $n$ pair \ill{} \ill{} $D_n$ \ill{}
$\neq \{1\}$}
\note{en haut à droite, un petit bloc de quatre lignes biffé d'obliques
(on y devine « réflexion », « $\mathbb{R}$ », « d'ordre 2 »), illisible
pour le reste.}

\page{45}
\drawing{un hexagone inscrit dans un cercle, de sommets marqués
$\tilde{R}_1$, $\tilde{R}'_0$, $\tilde{R}_5$, $\tilde{R}'_4$, $\tilde{R}_3$,
$\tilde{R}_2$ (ce dernier surchargé d'un $4$), le centre pointé ; les
diagonales et quelques cordes sont tracées en tirets au crayon.}

On s'intéresse au complexifié $V_{\mathbb{C}}$, et à la droite projective
complexe
\[
  \text{\struck{$\Sigma \ldots$}}
\]
\[
  (10)\qquad \Sigma_{\tilde{J}}(\mathbb{C}) = \Sigma_{\tilde{J}}
  \overset{\text{déf}}{=} \Sigma = \mathbb{P}(V_{\mathbb{C}})
  = V_{\mathbb{C}}^{*}/\mathbb{C}^{*} ,
\]
\add{les} \uncertain{sur lequel} $D$ opère \struck{\ill{}} au moyen de cette
opération ; le centre $\{1, a_0\}$ \uncertain{opérant trivialement}, \ill{}
$D/\{1, a_0\}$ (groupe diédral $D_n$). On \uncertain{note} $R_i \in \Sigma$
l'image de $\tilde{R}_i$ dans $\Sigma$, \ill{}
$R_i = R_{a_0(i)}$, \ill{} $J = \tilde{J}/\{1, a_0\}$ \ill{}.
\note{les lignes de ce paragraphe sont écrites en oblique, de bas à gauche
vers le haut à droite, entremêlées avec le NB qui suit ; seules les
formules se lisent avec quelque sûreté.}

\marginal{NB $(\Sigma, J \subset \Sigma)$ étant donnés ($\Sigma$ comme droite
projective complexe, (11) \ill{}) ainsi que \uncertain{sphère conforme
orientée} \ill{} $D_n$ \ill{} $(V_{\mathbb{R}}, \tilde{J})$ \ill{} à isom.\
près \ill{} $D_n$ !) \ill{}, mais seulement \ill{} les $D_n$ \ill{}
$\#J$ \ill{} torseurs \ill{} \ill{} modulo twist par \ill{} (11) \ill{} ;
\uncertain{l'ambiguïté} \ill{} à savoir \ill{} la donnée d'une \ill{} les
deux \ill{}-gônes dont les \ill{} des (13) autres \ill{}-gônes \ill{}}

\page{46}
\note{p.~3 de l'auteur.}
L'ens.\ des points fixes de $D^{+}$ opérant sur $\Sigma$ est
\struck{\uncertain{index}} réduit à deux éléments, un dans chacun des deux
hémisphères en lesquels $\Sigma_{\mathbb{R}} = P(V)$ découpe $\Sigma$ ---
on a :
\[
  (14)\qquad \pi_0(\Sigma \setminus \Sigma_{\mathbb{R}}) \simeq
  \underline{\omega} = \Omega(\Sigma_{\mathbb{R}})
\]
(\uncertain{couvr.} de Stokes, provenant de l'orientation can.\ de
$\Sigma$) donc les deux « pôles » seront indexés par $\underline{\omega}$,
et notés $P_\omega$
\[
  (15)\qquad
  \begin{cases}
    P_\omega \quad (\omega \in \underline{\omega}) \\
    P_\omega = \Sigma^{D^{+}} \cap \Sigma_\omega
  \end{cases}
\]
\note{sous $\Sigma^{D^{+}}$ : « invariants $D^{+}$ op.\ sur $\Sigma$ » ;
sous $\Sigma_\omega$ : « comp.\ conn.\ de $\Sigma \setminus \Sigma_{\mathbb{R}}$
indexée par $\omega$ ».}

Si on exprime la définition de $\Sigma^{D^{+}}$ en termes d'algèbre
linéaire, on trouve que ce sont les droites \add{propres} (dans $V_{\mathbb{C}}$)
\struck{propres} pour l'action de $D^{+}$. Mais considérons la
\uncertain{double} structure euclidienne sur $V_{\mathbb{C}}$, qui
\struck{définit} s'étend en une structure

\page{47}
\uncertain{quadratique} non dégénérée sur $V_{\mathbb{C}}$, par laquelle
$V_{\mathbb{C}}$ se décompose en sous-espaces isotropes
\[
  (16)\qquad V_{\mathbb{C}} = \bigoplus_{\omega \in \underline{\omega}}
  V_{\mathbb{C}}^{\omega}
\]
qui sont aussi les sous-espaces propres pour le groupe \add{$O^{+}(V)$} des
rotations de $V$. La donnée d'une orientation \add{$\omega$} de $V$ permet de
choisir une \uncertain{autom.}\ de
\[
  (17)\qquad
  \begin{cases}
    u_\omega \in \text{\struck{\ill{}}}\; O^{+}(V) \\
    u_\omega^{2} = -1
  \end{cases}
\]
\marginal{rotation d'un quart de tour dans le sens de l'orientation
$\omega$}
d'où une structure complexe sur $V$ \add{(qui devient un vect.\ complexe de
rang 1, soit $V_{(\omega)}$)}, en \struck{faisant} opérer $i \in \mathbb{C}$
par
\[
  i_\omega\, x = u_\omega x .
\]
Mais par la définition universelle du complexifié, la \struck{structure}
application $\mathbb{R}$-linéaire $V \to V_{(\omega)}$ se prolonge
canoniquement en une application

\page{48}
\note{p.~4 de l'auteur.}
$\mathbb{C}$-linéaire, \struck{commutant d'ailleurs (par transport de str.)}
\[
  (17')\qquad V_{\mathbb{C}} \xrightarrow{\;p_\omega\;} V_{(\omega)}
  \qquad (\mathbb{C}\text{-lin.})
\]
\marginal{cette application commute par transport de str.\ à $O^{+}(V)$
(opérant sur $V_{\mathbb{C}}$ par complexification)}
d'où une application \add{$\mathbb{C}$-lin.} canonique
$V_{\mathbb{C}} \to \prod_{\omega \in \underline{\omega}} V_{(\omega)}$,
qui est un isomorphisme
\[
  (18)\qquad p = (p_\omega)_{\omega \in \underline{\omega}} : V_{\mathbb{C}}
  \xrightarrow{\;\sim\;} \prod_{\omega \in \underline{\omega}}
  \underbrace{V_{(\omega)}}_{\simeq V_{\mathbb{C}}^{\omega}} .
\]
\marginal{\uncertain{l'opération} d'\ill{} $V_{\mathbb{C}}$ \ill{}
$V^{(\omega)}$ \ill{} comme \ill{} par \ill{} $u_\omega$ sur
\uncertain{le facteur}}
Les $V_{(\omega)}$ sont apparus comme espaces quotients de $V_{\mathbb{C}}$,
si on veut les regarder comme sous-espaces \add{--- auquel cas on les note
$V_{\mathbb{C}}^{\omega}$ ---} on trouve que
\[
  (19)\qquad
  \begin{aligned}
    V_{\mathbb{C}}^{\omega}
    &= \operatorname{Ker}\bigl(V_{\mathbb{C}} \xrightarrow{\;p_{\omega'}\;}
       V_{(\omega')}\bigr) \\
    &= \{\, x \in V_{\mathbb{C}} \mid u_{\omega,\mathbb{C}}(x) = i.x \,\}
       \qquad (\omega \in \underline{\omega})
  \end{aligned}
\]
\marginal{Aussi $V_{\mathbb{C}}^{\omega} = \{\, x - i u_\omega x \mid x \in V \,\}$,
\ill{} $\varphi_\omega : V \xrightarrow{\sim} V_{\mathbb{C}}^{\omega}$ le
$\mathbb{R}$-isom.\ induit par (17')}
En fait, les $V_{\mathbb{C}}^{\omega}$ \add{sont des sous-espaces} propres
pour l'action de $O^{+}(V)$, i.e.\ stables par cette action, tandis que les
éléments de $O^{-}(V)$ les interchangent.
\marginal{identification \ill{} des $V^{\omega}$ \ill{} propres \ill{}
\uncertain{coïncide} \ill{} par \ill{} Stokes}

On s'intéresse au fibré vectoriel

\page{49}
$\mathcal{O}_\Sigma(1)$ sur $\Sigma$, défini par la représentation (10) de
$\Sigma$ comme \struck{fibré} \uncertain{espace des} droites projectives
associé à un vectoriel bien déterminé. Je désigne par
$P\Sigma_{J,\mathbb{C}}$ --- \uncertain{simplement} $P\Sigma$ --- le fibré
principal associé, de groupe $\mathbb{C}^{*}$ (avec la topologie complexe,
disons, avec la topologie habituelle --- pour le moment \ill{} dans le
contexte des schémas…). \struck{\ill{}}

\uncertain{Un} point $x$ de $\Sigma$ étant une droite $F$ de $V$,
\struck{\ill{}} on a des isom.\ can.
\[
  (18)\qquad
  \begin{cases}
    V(\underline{\mathcal{O}}(1))_x \simeq V/F
      \qquad \bigl(\mathcal{O}_F(1) \overset{\text{déf}}{=} \ldots\bigr) \\
    F \otimes \mathcal{O}_F(1) \simeq \det(V) \simeq
      \mathbb{R}(\underline{\omega})
      \;\bigl(\overset{\text{déf}}{=} \mathbb{R} \wedge_{\{\pm1\}}
      \underline{\omega}\bigr) \\
    \mathcal{O}_F(1) \simeq \check{F}(\omega)
      \;\bigl(\overset{\text{déf}}{=} \check{F} \wedge_{\{\pm1\}}
      \underline{\omega}\bigr) \simeq F^{\perp}
  \end{cases}
\]
\note{le numéro (18) est repris tel quel, bien qu'il serve déjà
p.~48 ; de même (19) ci-dessous. En marge de la première ligne,
« d'où, pour $x$ : dét » ; au-dessus de la troisième, « $\mathbb{R}$-tor. ».}
\marginal{orth.\ pour la str.\ quadratique \uncertain{canonique}}
d'où
\[
  (19)\qquad (P\Sigma)_x = (P\Sigma)_F \simeq (V/F)^{*} \simeq
  (\check{F})^{*} \wedge_{\{\pm1\}} \underline{\omega} \simeq F^{\perp}
\]
\struck{et par son} isom.\uncertain{s} compatibles avec l'action de
$\mathbb{C}^{*}$ (opérant sur $\check{F}$ par \struck{contra}
\uncertain{contragrédience} \uncertain{de l'action} sur $F$), on a encore
\[
  (20)\qquad (P\Sigma)_x \simeq F^{*} \wedge_{\{\pm1\}} \underline{\omega}
  = (F \wedge_{\{\pm1\}} \underline{\omega})^{*} \simeq (F^{\perp})^{*}
\]

\page{50}
\note{p.~5 de l'auteur.}
\ill{} un isom.\ (\ill{} $V(\omega)^{*}$, $V^{*}$) pour deux droites
vectorielles $F, F'$ \ill{} $F^{*} \simeq F'^{*}$ \ill{}, entre
$(P\Sigma)_x$ et $(F \wedge \omega)^{*}$, \struck{ind}
\uncertain{qui renverse} l'action \ill{}. On trouve ainsi des isom.
\[
  (21)\qquad
  \begin{cases}
    P\Sigma \simeq (V(\omega))^{*} \\
    P\Sigma \simeq V^{*}
  \end{cases}
\]
\note{à droite de (21), quelques symboles biffés (« $F$ », « $\mapsto
\check{e}$ », « $F/e$ »).}
\marginal{\ill{} $(P\Sigma)_x$ \ill{} $\simeq \check{F}(\omega) \simeq
F^{\perp}$ \ill{} $(F^{\perp})^{*}$ \ill{} (21) \ill{} $D$ \ill{}
l'action de \ill{} (\uncertain{transformant} \ill{})}
Mais on a un isom.\ canonique (provenant du produit extérieur)
\[
  V \simeq \check{V}(\omega)
\]
et un isom.\ canonique (provenant de la forme quadratique)
$V \to \check{V}$, donc $\check{V}(\omega) \simeq V(\omega)$, d'où un
isom.\ canonique
\[
  (21)\qquad V \xrightarrow{\;\sim\;} V(\omega)
\]
\uncertain{qui explicite} la \ill{} de l'un à l'autre isom.\
\struck{\ill{} (21)}. D'ailleurs, (21) \uncertain{induit}
\struck{un isom.\ $D$, et par \ill{} induisant l'identité}
\[
  \text{\struck{$\Sigma \simeq P(V) \xrightarrow{\text{id}}
  \Sigma \simeq P(V(\omega))$}}
\]
\[
  (22)\qquad \Sigma = P(V) \xrightarrow{\;\sigma\;} P(V(\omega))
  \qquad (\simeq \Sigma = P(V))
\]
(\struck{qui est} la \struck{rotation} symétrie p.r.\ à l'axe polaire,
\struck{qui \ill{}} qui s'exprime \struck{\ill{}} comme \uncertain{le}
passage de $F$ à $F^{\perp}$).

\struck{\ill{} est iso \ill{} si $F$ correspond à $F' \subset V$, on fait
$F'(\omega)$ \ill{} dans $V(\omega)$ \ill{} $F' = F^{\perp}$) on a
$F(\omega) \simeq F'(\omega)$ i.e.\ $F(\omega) \simeq F'$.}
\note{ces trois dernières lignes sont barrées en bloc de traits obliques.}

\page{51}
\struck{\ill{}} \ill{} l'iso.
\[
  V \xrightarrow{\;\sim\;} \ldots \Sigma \simeq V^{*}
\]
\[
  P\Sigma \simeq V(\underline{\omega})^{*} ,
\]
\struck{\ill{}} \add{\ill{}} l'action de $\mathbb{C}^{*}$ \struck{\ill{}}
(\ill{} par l'action opposée), \ill{} compatible \ill{} \ill{}
projections
\begin{tikzcd}
  P\Sigma \arrow[r, "\sim", "(23)"'] \arrow[d, "\text{proj.\ can.}"']
    & V^{*} \arrow[d, "\text{proj.\ can.}"] \\
  \Sigma \arrow[r, "\sigma"] & \Sigma \simeq V^{*}/\mathbb{R}^{*}
\end{tikzcd}

Notons en passant que si $\tau$ désigne la conjugaison complexe (dans
$V_{\mathbb{C}}$ --- \ill{} $\Sigma$) on a
\[
  \tau\sigma = \underline{a}
\]
« antipodisme » de $\Sigma$…
\note{le numéro (25) figure dans la marge à hauteur de cette formule ;
(24) n'apparaît que dans la note marginale.}

Je m'intéresse à l'ouvert
\[
  (26)\qquad \Sigma^{*} = \Sigma \setminus \{\, R_i \mid i \in J \,\}
\]
et à
\[
  (27)\qquad P\Sigma^{*} = P\Sigma \,|\, \Sigma^{*}
  \overset{(23)}{\simeq} V^{**} = V^{*} - \bigcup_{i \in J} \Delta_i
\]
où les $\Delta_i$ sont les $n = m/2$ \emph{diagonales} du
\marginal{Tout le \ill{} digression précédente \ill{} \uncertain{« sens des}
\ill{} : \ill{} d'un point $R_i \in \Sigma$, il \ill{} $R_i$ canonique
(\ill{} $\omega \in \underline{\omega}$), et \ill{}
\struck{$\Gamma(\Sigma, L)$} $\Gamma(\Sigma, \underline{L})$ \ill{} $R_i$
des éléments \ill{} de $\Gamma(\Sigma, L)$ \ill{} dont la somme est $0$,
\ill{} $V \simeq \mathcal{O}^{J}/\mathcal{O}$ (diagonale) \ill{} $\Sigma$
\ill{} $V^{*}/\mathbb{C}^{*}$ \ill{} $\Sigma$ : \ill{} $P\Sigma$ \ill{}
$\mathbb{C}^{*}$ \ill{} l'opération de \ill{} (24) $V = \Gamma(\Sigma,
\underline{L})(\omega)$ \ill{} $V = \operatorname{Ker}(\mathcal{O}^{J} \to
\mathcal{O})$ \ill{} $D\ (\simeq \mathfrak{S}_3)$ \ill{} (25)}
\note{la note marginale ci-dessus est écrite en oblique sur tout le quart
gauche de la feuille, en travers d'elle-même ; seuls des fragments s'en
lisent.}

\page{52}
\note{p.~6 de l'auteur.}
polygone régulier. Je veux décrire des groupoïdes équivalents aux
groupoïdes fondamentaux
\[
  (28)\qquad \Pi_1 P\Sigma^{*} \simeq \Pi_1 V^{**} , \qquad \Pi_1 \Sigma^{*}
\]
par générateurs et relations, avec des sommets bien choisis dans les
espaces envisagés, stables par l'action de $D$ ainsi que de celle de
$\tau$. Bien sûr, il faut garder présent le morphisme de groupoïdes
\[
  (29)\qquad \Pi_1 P\Sigma^{*} = \Pi_1 V^{**} \longrightarrow \Pi_1 \Sigma^{*} .
\]
Sur la situation (28), (29) opère le groupe
\[
  \text{\struck{\ill{}}}\; D \times \mathbb{C}^{*} \times \{1, \tau\}
\]
\marginal{$D$ opère \ill{}, $\mathbb{C}^{*}$ par homothéties (mais c'est
l'action opposée de celle sur $P\Sigma^{*}$)}
\marginal{NB $D$ n'opère sur $\Sigma^{*}$, $P\Sigma^{*}$ que \ill{} \ill{}
(\ill{} \uncertain{pas dans le cas} $n = 3$) \struck{\ill{}}}
\uncertain{NB} action \uncertain{non fidèle}, l'élément $(a_0, -1)$
\uncertain{opérant trivialement}.

On s'intéresse \add{fini} à l'action des sous-groupes
\[
  D \times \mu_m \times \{1, \tau\}
\]
d'ordre \struck{\ill{}} $4m^{2} = 16n^{2}$ ($144$ si $n = 3$) mais
\uncertain{qui n'est} pas fidèle --- le quotient \uncertain{qui opère}
fidèlement est d'ordre $8n^{2}$ ($72$ si $n = 3$).
\marginal{NB l'action \uncertain{quotient} (30*)
$(D \times \mu_m)/\{\pm 1\} \times \{1, \tau\}$ \uncertain{qui opère
fidèlement}}
\note{le numéro (30) est porté dans la marge, à côté de la note
ci-dessus ; il semble désigner le groupe $D \times \mu_m \times \{1,\tau\}$.}
On peut \uncertain{s'intéresser} plutôt à son sous-groupe
\[
  (31)\qquad D^{+} \times \mu_m \simeq \mu_m(\underline{\omega}) \times \mu_m
\]
de (30), en notant que pour la représentation

\page{53}
(16), (18), i.e.\ $V_{\mathbb{C}} \simeq \bigoplus_{\omega \in
\underline{\omega}} V_{\mathbb{C}}^{\omega}$, \uncertain{on a} le groupe
produit $\mu_m^{\underline{\omega}}$ d'ordre $m^{2}$ opère sur
$V_{\mathbb{C}}$ : par
\[
  (32)\qquad (\lambda_\omega)_{\omega \in \underline{\omega}} \cdot
  (x_\omega)_{\omega \in \underline{\omega}} = (\lambda_\omega x_\omega)
\]
On \struck{va} a des homomorphismes
\[
  (33)\qquad
  \begin{cases}
    \mu_m \longrightarrow \mu_m^{\underline{\omega}}
      \qquad \text{hom.\ diagonal} \\
    \mu_m(\underline{\omega}) \simeq D^{+} \longrightarrow
      \mu_m^{\underline{\omega}} \\
    \qquad \lambda \longmapsto (\lambda \wedge \omega)_{\omega \in
      \underline{\omega}}
  \end{cases}
\]
\note{à droite de la dernière ligne, un mot biffé et, au-dessus,
\uncertain{« droit »}.}
tels que les actions de $\mu_m$, $D^{+}$ sur $V^{**}$ \struck{\ill{}} se
font via celle de $\mu_m^{\underline{\omega}}$. L'hom.\ correspondant
\[
  (34)\qquad \text{\struck{\ill{}}}\;
  \underbrace{D^{+}}_{\mu_m(\underline{\omega})} \times \mu_m
  \longrightarrow \mu_m^{\underline{\omega}}
\]
\uncertain{entre} groupes de même ordre $m^{2} = 4n^{2}$ a comme noyau le
sous-groupe d'ordre 2 engendré par \struck{\ill{}} $(a_0, -1)$, son image est
donc d'indice 2, elle est formée des couples
$(\lambda_\omega)_{\omega \in \underline{\omega}}$ tels que l'on ait
\[
  (35)\qquad \text{\struck{\ill{}}}\;
  \Bigl(\prod_{\omega \in \underline{\omega}} \lambda_\omega\Bigr)^{n} = 1
  \qquad \text{i.e.}\quad (\lambda_\omega \lambda_{\omega'})^{n} = 1 .
\]
Notons que les autom.\ de $V_{\mathbb{C}}$ qui stabilisent
$V_{\mathbb{C}}^{**}$ i.e.\ qui induisent un \uncertain{automorphisme} de
$\Sigma$ qui stabilise $J$, sont ceux qui provien-

\page{54}
\note{p.~7 de l'auteur.}
nent de $D$ mod $\mathbb{C}^{*}$, donc qui proviennent de
$D \times \mathbb{C}^{*}$. Ceux qui conservent \struck{\ill{}} chaque
$P_\omega$, i.e.\ chaque sous-espace isotrope $V_{\mathbb{C}}^{\omega}$,
sont ceux qui proviennent de $D^{+} \times \mathbb{C}^{*}$. Si on plonge
\[
  (36)\qquad D^{+} \times \mathbb{C}^{*} \hookrightarrow
  \mathbb{C}^{*\,\underline{\omega}}
\]
comme dans (34), alors l'image est formée des éléments
$(\lambda_\omega)_{\omega \in \underline{\omega}}$ tels que, choisissant
$\omega \in \underline{\omega}$, on ait
\[
  (37)\qquad \text{\struck{\ill{}}}\;
  (\lambda_\omega \lambda_{\omega'}^{-1})^{n} = 1
\]
(condition indépendante du choix de $\omega$), ce qui, pour
\struck{\ill{}} $\lambda_\omega, \lambda_{\omega'} \in \mu_m = \mu_{2n}$,
équivaut à la condition sous la forme (35). Donc plus symétriquement,
\uncertain{les} éléments de $\mu_m^{\underline{\omega}}$ qui ne
satisfont pas cette relation ne stabilisent pas $V^{**}$.

On s'intéresse à des choix de pts base \ill{}, un \struck{\ill{}}
\uncertain{état de cases}, \uncertain{suivant} stables par
$D \times \{1, \tau\}$, et éventuellement --- par le groupe de symétries
plus grand
\[
  (D \times \mu_m)/\{\pm 1\} \times \{1, \tau\}
\]
(d'ordre 72 dans le cas $n = 3$)
\note{le texte rédigé (pp.~42--54, p.~1 à 7 de l'auteur) s'arrête ici, en
bas de feuille ; la suite n'est pas dans ce lot.}

\page{55}
${}_0\mathfrak{T}_{0,5}$ et $S_0\mathfrak{T}_{0,5}$
\note{seule ligne de la feuille, au crayon, en haut à droite : titre de
ce qui suit. La p.~56 est blanche.}

\page{57}
\drawing{en haut à gauche, une sphère portant un grand cercle (équateur)
avec quatre petits cercles pointés ; légende « 4 sommets, deux faces ».
Une flèche « deux façons » mène à une seconde sphère, même équateur à quatre
sommets cerclés, coupé par un grand cercle méridien (« cercle marqué ») ;
légende « 1 cercle marqué, 6 sommets, 4 faces ».}

3 possibilités \uncertain{sur} I (trois faces aux centres \uncertain{gauches})
; 6 possibilités sur I (6 arêtes du cercle \uncertain{gauche})

\drawing{une flèche vers le haut, marquée « trois façons (supprimer une des
trois \struck{\ill{}} d'arêtes opposées) », part d'un tétraèdre dont les
quatre sommets portent de petits angles ; légende « 4 sommets, 4 faces ».
À droite, une petite sphère au crayon et un croquis raturé (une sphère
quadrillée dans un cube), avec « \struck{\ill{}} I ».}

deux possibilités sur I

\drawing{sous un trait horizontal, deux sphères. À gauche, « I », avec le
schéma de trois branches « $\vee\!\!\!\perp\!\vee$ » : équateur à cinq
sommets cerclés, coupé par deux grands cercles méridiens ; « $\mathbb{Z}/2$ » ;
légende « $9 = 5 + 4$ sommets, \struck{\ill{}} faces, deux cercles marqués »
et « $12.5 = 60$ possibilités pour $K$ (pentagone + sommet) », « $D_{01}$ ».
À droite, « II », schéma « $\vee\!\!\perp\!\vee$ » : équateur à cinq sommets
cerclés, un méridien ; « $\mathbb{Z}/2$ » ; légende « $5 + 2 = 7$ sommets,
4 faces, 1 cercle marqué » et « $12.5 = 60$ poss.\ \uncertain{pour}
\ill{} (pentagone + arête) », « $D_{02}$ ».}
\note{« $12.5 = 60$ » sous II est écrit « 66 » ou « 60 » ; le calcul
$12 \cdot 5$ donne 60.}


\page{59}
\drawing{à gauche, « III » : une sphère portant deux cercles parallèles
reliés par deux montants verticaux, le cercle du haut « cercle marqué »,
deux sommets cerclés sur le cercle du bas et deux au pôle ; « $\mathbb{Z}/2$ » ;
légende « $9 = 3 + 3 + 1 + 2$ sommets, 6 faces, 1 cercle marqué » ;
« $60 = 10 \times 3 \times 2$ poss.\ pour $K$ » ;
« (triangle + sommet marqué + \ill{}) » avec sous les termes 10, 3, 2.
Au centre, un schéma en « $\vee\!\!\perp\!\vee$ » à deux branches. À
droite : « $D_0$ », « III » : une sphère à équateur portant cinq sommets
cerclés, marquée « $D_5$ » ; légende « \struck{12 façons} 5 sommets, 2 faces,
12 possibilités qui \ill{} orientation de $K$ » ; « IV », « $D_2$ » : une
sphère méridienne en fuseaux avec trois sommets cerclés ; légende
« \struck{\ill{}} sommets, 6 faces, 20 possibilités (cf \ill{}, car il
faut choisir une orientation de la sphère) $\sim$ orientation de $K$ » ;
un croquis raturé « $\mathfrak{S}_3$ ».}

choisissons une orientation sur $K$, \ill{} et \ill{}, d'abord dans
I, II, \uncertain{avec} \uncertain{pentagones} compatibles \ill{} l'orienta-
tion. \struck{$\mathrm{II} \to \mathrm{I}$ deux façons} Alors
\begin{itemize}
  \item donnée I $=$ sommet de l'icosaèdre \add{$\mathrm{Ic}_\omega$}
    \add{\uncertain{grandes}} $+$ arête \struck{\ill{}} \add{incident},
    i.e.\ $\in D_{01}(\mathrm{Ic}_\omega)$ \uncertain{diagramme} de type 01 ;
  \item donnée II $=$ sommet de l'icos.\ \ill{} $+$ face \struck{\ill{}}
    \add{incident}, i.e.\ $\in D_{02}(\mathrm{Ic}_\omega)$ ;
  \item donnée III $=$ \struck{\ill{}} \add{\uncertain{resp.} de} \ill{} de
    l'icos.\ \struck{\ill{} + sommet dessus}
\end{itemize}

\drawing{en bas, plusieurs croquis largement raturés : une ellipse, une
sphère en fuseaux à quatre sommets cerclés (« $9 = 5 + 4\,\struck{+1}$
sommets », « 30 possibilités (car il faut choisir l'orientation du carré) »),
une petite sphère à quatre sommets, un tétraèdre, une sphère « $D_1$ » ;
flèches « $\to D_1$ », « $\mathbb{Z}$ ». À gauche, une liste :}
\[
  \begin{array}{l}
    \text{\struck{$D_0 \to D_1$ \quad cinq façons}} \\
    D_1 \to D_2 \\
    D_{01} \to D_0 \\
    D_{01} \to D_1 \\
    D_{02} \to D_0 \\
    D_{02} \to D_2 \quad \text{\struck{\ill{}}}
  \end{array}
\]


\end{document}
